% Part: lambda-calculus
% Chapter: introduction
% Section: unique-readability

\documentclass[../../../include/open-logic-section]{subfiles}

\begin{document}

\olfileid{lam}{syn}{unq}
\olsection{एकमेव वाचनीयता}

प्रत्येक पद घडवण्याचा एकमेव मार्ग असतो का, असा प्रश्न
पडू शकतो; आणि तसेच आहे. प्रत्येक लॅम्डा पद घडवण्याची
आणि त्याचा अर्थ लावण्याची एकच पद्धत आहे. पुढील चर्चेत
\emph{रचनाप्रक्रिया} म्हणजे \olref[trm]{defn:term}
मधील रचनानियम एकदा किंवा अनेकदा वापरून पद घडवण्याची
प्रक्रिया.

\begin{lem}\ollabel{lem:term-start}
पदाची सुरुवात एकतर चराने किंवा कंसाने होते.
\end{lem}

\begin{proof}
\olref[trm]{defn:term} नुसार घडवले असेल तरच एखादी
गोष्ट पद ठरते. ते \olref[trm]{defn:term-var} चे फलित
असेल, तर ते चर असलेच पाहिजे. ते
\olref[trm]{defn:term-abs} किंवा
\olref[trm]{defn:term-app} चे फलित असेल, तर
त्याची सुरुवात कंसाने होते.
\end{proof}

\begin{lem}\ollabel{lem:app-start}
उपयोजनाच्या फलिताची सुरुवात एकतर दोन कंसांनी किंवा
एक कंस आणि त्यानंतर चर याने होते.
\end{lem}

\begin{proof}
$M$ हे उपयोजनाचे फलित असेल, तर ते $(PQ)$ या रूपात
असते; म्हणून त्याची सुरुवात कंसाने होते. $P$ हे पद
असल्यामुळे \olref{lem:term-start} नुसार त्याची
सुरुवात कंसाने किंवा चराने होते.
\end{proof}

\begin{lem}\ollabel{lem:initial}
पदाचा कोणताही उचित आरंभीचा भाग स्वतः पद नसतो.
\end{lem}

\begin{prob}
पदांच्या लांबीवरील विगमनाने
\olref[lam][syn][unq]{lem:initial} सिद्ध करा.
\end{prob}

\begin{prop}[एकमेव वाचनीयता] \ollabel{prop:unq}
प्रत्येक पदाची रचनाप्रक्रिया एकमेव असते. दुसऱ्या
शब्दांत, एखाद्या रचनाप्रक्रियेने पद $M$ घडत असेल,
तर तेच पद घडवू शकणारी ती एकमेव रचनाप्रक्रिया असते.
\end{prop}

\begin{proof}
पदांच्या रचनाप्रक्रियेवरील विगमनाने हे सिद्ध करू.
  \begin{enumerate}
    \item $M$ हे $x$ या रूपात आहे, जिथे $x$ हे
      एखादे चर आहे. अमूर्तीकरण आणि उपयोजन यांची
      फलिते नेहमी कंसाने सुरू होतात, म्हणून $M$
      घडवताना त्यांचा वापर झालेला असू शकत नाही.
      त्यामुळे $M$ ची रचनाप्रक्रिया
      \olref[trm]{defn:term}\olref[trm]{defn:term-var}
      ची एकच पायरी असली पाहिजे.
    \item $M$ हे $(\lambd[x][N])$ या रूपात आहे,
      जिथे $x$ हे चर आणि $N$ हे पद आहे. ते
      \olref[trm]{defn:term}\olref[trm]{defn:term-var}
      नुसार घडलेले असू शकत नाही, कारण ते एकटे
      चर नाही. \olref{lem:app-start} नुसार ते
      उपयोजनाचेही फलित नाही. त्यामुळे $M$ हे
      फक्त~$N$ वरील अमूर्तीकरणाचे फलित असू शकते.
      विगमन गृहीतकानुसार $N$ ची रचनाप्रक्रियाही
      एकमेव आहे.
    \item $M$ हे $(PQ)$ या रूपात आहे, जिथे $P$
      आणि $Q$ ही पदे आहेत. त्याची सुरुवात कंसाने
      होत असल्याने ते
      \olref[trm]{defn:term}\olref[trm]{defn:term-var}
      नुसारही घडलेले असू शकत नाही.
      \olref{lem:term-start} नुसार $P$ ची सुरुवात
      $\lambd$ ने होऊ शकत नाही; म्हणून $(PQ)$ हे
      अमूर्तीकरणाचे फलितही असू शकत नाही. आता $M$
      उपयोजनाने घडवण्याचा आणखी एक मार्ग आहे असे
      समजा; उदाहरणार्थ, ते $(P'Q')$ या रूपातही
      आहे. सुरुवातीची पदे वेगळी असतील, तर $P$
      हा $P'$ चा उचित आरंभीचा भाग असेल किंवा
      उलट असेल; \olref{lem:initial} नुसार हे
      अशक्य आहे. ती एकच असतील, तर उरलेले भागही
      एकच असतात. म्हणून $P$ आणि $Q$ एकमेव
      रीतीने ठरतात आणि विगमन गृहीतकानुसार $P$
      आणि $Q$ यांच्या रचनाप्रक्रियाही एकमेव असतात.
  \end{enumerate}
\end{proof}

वरील प्रतिज्ञेचा अधिक वाचनीय अर्थ असा आहे:
\begin{prop}
  पद $M$ पुढीलपैकी एकाच रूपात असू शकते:
  \begin{enumerate}
    \item $x$, जिथे $x$ हे $M$ ने एकमेव रीतीने
      ठरणारे चर आहे.
    \item $(\lambd[x][N])$, जिथे $x$ हे चर आणि
      $N$ हे दुसरे पद आहे; दोन्ही $M$ ने एकमेव
      रीतीने ठरतात.
    \item $(PQ)$, जिथे $P$ आणि $Q$ ही दोन पदे
      $M$ ने एकमेव रीतीने ठरतात.
  \end{enumerate}
\end{prop}

\end{document}
